\subsection{Use of Generative Artificial Intelligence Tools}

Generative Artificial Intelligence (AI) tools were employed to support various stages of the documentation and design process in this project. The tools used include \textbf{Perplexity AI} and \textbf{Gemini}, selected for their advanced natural language generation and reasoning capabilities. These tools assisted in tasks related to technical writing, software modeling, and language revision, always under careful human supervision and validation.

\textbf{All strategic and main decision processes were entirely human-driven}. AI served strictly as a tool to augment these decisions, automate routine work, and save time, without influencing core architectural or methodological choices.

The specific prompts provided to the AI systems were explicitly formulated to focus on technical and linguistic support. Examples of inputs include: 
\begin{itemize}
    \item ``Help me generate a PlantUML code Domain Class Diagram from the following Domain Class Diagram sketch.''
    \item ``Help me brainstorm goals for the following software use case.''
    \item ``Help me write in a formal way the following functional requirements.''
    \item ``Help me rewrite the following section for a software design document.''
    \item ``Help me look for grammar and spelling mistakes.''
\end{itemize}

The outputs obtained included automatically generated textual content such as structured PlantUML code for all diagrams, rewritten and polished technical documentation, formalized versions of functional requirements, and identified language corrections. In brainstorming tasks, the tools provided diverse options and design considerations that served as inspiration for choosing what to include in the document.

Verification, refinement, and integration of AI-generated outputs were performed systematically. Each generated segment was carefully reviewed before inclusion. For code artifacts (Alloy), outputs were tested directly within the target environment or modeling software to confirm proper functionality and consistency with the project's requirements. When AI responses incorporated external knowledge sources, these references were verified for accuracy and reliability before use. Sections of text that did not meet quality or factual standards were either revised or discarded entirely.

The integration of AI assistance enhanced productivity and coherence across the project documentation. It supported the generation of consistent technical descriptions and facilitated the refinement of complex technical ideas into well-structured prose. Human oversight remained indispensable; all final content (including diagrams, requirements, and textual documentation) was critically validated, edited, and approved by the authors to ensure academic rigor, correctness, and originality.

Overall, Generative AI tools functioned as an auxiliary mechanism for idea development, code drafting, and language improvement, rather than as a replacement for human analysis or authorship. Through responsible use and verification, these tools contributed to improving the clarity, precision, and quality of the final project deliverables.
